\begin{note}
	Далее рассматриваем $d^2_{\vx\vx}L$, в котором $\v\la$ фиксируются, а $dx_i$ подчиняются продифференцированным уравнениям связи $d\phi_i(\vx)=0,\ i=1,\dots,m$, то есть при дополнительном условии $\ds\frac{\mc D(\phi_1,\dots,\phi_m)}{\mc D(x_1,\dots,x_m)}(\vx\n)\ne0$ $d^2_{\vx\vx}L$ является квадратичной формой $dx_{m+1},\dots,dx_{n}$.
\end{note}
\begin{anote}
	Ещё раз напишу эту матрицу, чтобы гиперссылка вела не так далеко:
	\begin{equation}\label{29_m1}
		\begin{pmatrix}
			\frac{\vd \phi_1}{\vd x_1} & \dots & \frac{\vd \phi_1}{\vd x_n} \\
			\vdots & \ddots & \vdots \\
			\frac{\vd \phi_n}{\vd x_1} & \cdots & \frac{\vd \phi_n}{\vd x_n}
		\end{pmatrix}
	\end{equation}
\end{anote}
\begin{theorem}\nf{(Достаточное условие условного экстремума)}\label{29_th4}\\
	Если $f,\phi_1,\dots,\phi_m$ дважды непрерывно дифференцируемы в окрестности точки $\vx\n$, матрица Якоби (\ref{29_m1}) имеет ранг m в точке $\vx\n$, $(\vx\n,\v\la\n)$ является стационарной точкой функции Лагранжа $L(\vx,\v\la)=f(\vx)+\sum_{i=1}^m\la_i\phi_i(\vx)$, то:
	\begin{enumerate}
		\item В случае положительной (отрицательной) определённости $d^2L_{\vx\vx}(\vx\n)$, как квадратичной формы переменных $dx_1,\dots,dx_n$ подчинённых продифференцированным уравнениям связи $d\phi_i(\vx\n)=0,\ i=1,\dots,m$, точка $\vx\n$ является точкой условного минимума (максимума) функции $f$ при наличии связей $\phi_i(x_1,\dots,x_n)=0,\ i=1,\dots,m$.
		\item В случае знаконеопределённости указанной формы $d^2_{\vx\vx}L(\vx\n)$ точка $\vx\n$ не является точкой условного экстремума f при наличии связей $\phi_i(x_1,\dots,x_n)=0,\ i=1,\dots,m$.
	\end{enumerate}
\end{theorem}
\begin{proof}
	Без ограничений общности считаем, что $\ds\frac{\mc D(\phi_1,\dots,\phi_m)}{\mc D(x_1,\dots,x_m)}(\vx\n)\ne0$, Тогда задача на условный экстремум равносильна задаче на локальный (безусловный) экстремум функции $F(\vx)=f(\psi_1(\tvx),\dots,\psi_m(\tvx),\tvx)$, $\psi$ из $x_i=\psi_i(x_{m+1},\dots,x_n),\ i=1,\dots,m$.\\
	Достаточное условие локального экстремума функции F в точке $\tvx\n$ состоит в знакоопределённости $d^2_{\vx\vx}F(\tvx\n)$.\\
	Вычислим этот дифференциал с помощью инвариантности формы первого дифференциала:
	\[dF(\tvx)=df(\psi_1(\tvx),\dots,\psi_m(\tvx),\tvx)=df(\vx)=dL(\vx,\v\la)=\sum_{i=1}^n\frac{\vd L}{\vd x_i}(\vx,\v\la)dx_i\]
	Ранее доказывалось, что $d\psi_i(\vx)=0\lra dx_i=d\psi_i,\ i=1,\dots,m$, тогда вычислим второй:
	\begin{multline*}
		d^2F(\tvx)=\sum_{i=1}^nd\l(\frac{\vd L}{\vd x_i}\r)(\vx\n,\v\la\n)dx_i+\sum_{i=1}^n\frac{\vd L}{\vd x_i}(\vx\n,\v\la\n)d^2x_i=\\
		=\sum_{i,j=1}^n\frac{\vd^2L}{\vd x_i\vd x_j}(\vx\n,\v\la\n)dx_idx_j=d^2_{\vx\vx}L(\vx\n,\v\la\n)
	\end{multline*}
	Получили требуемое.
\end{proof}
\begin{example}
	Номер $5.32(2)$.\\
	Найти наименьшее и наибольшее значение функции $u=x^2+2y^2+3z^2$ при выполнении условий $x^2+y^2+z^2=1,\ x+2y+3z=0$.\\
	Для начала найдём локальные экстремумы.\\
	Функция Лагранжа: $L(x,y,z,\la,\mu)=x^2+2y^2+3z^2+\la(x^2+y^2+z^2-1)+\mu(x+2y+3z)$
	\begin{equation*}
		\begin{split}
			\frac{\vd L}{\vd x}&=2x+2\la x+\mu=0\\
			\frac{\vd L}{\vd y}&=4y+2\la y+2\mu=0\\
			\frac{\vd L}{\vd z}&=6z+2\la z+3\mu=0\\
			\frac{\vd L}{\vd\la}&=x^2+y^2+z^2=0\\
			\frac{\vd L}{\vd\mu}&=x+2y+3z=0
		\end{split}
	\end{equation*}
	\begin{equation*}
		\begin{cases}
			2+\frac{5\la}3x-\frac\la3y-\frac\la3z=0\\
			-\frac{2\la}3x+\l(4+\frac{4\la}3\r)y-\frac{2\la}3z=0\\
			-\la x-\la y+(6+\la)z=0
		\end{cases}
	\end{equation*}
	\begin{equation*}
		A=
		\begin{pmatrix}
			2+\frac{5\la}3 & -\frac\la3 & -\frac\la3\\
			-\frac{2\la}3 & 4+\frac{4\la}3 & -\frac{2\la}3\\
			-\la & -\la & 6+\la
		\end{pmatrix}
	\end{equation*}
	Теперь решим $|A|=0$. Решением является $\frac73\la^2+8\la+6=0$, то есть:
	\[\la_{1,2}=\frac{-4\pm\sqrt2}{7/3}\]
	$\ds U_{max}=\frac{3(4+\sqrt2)}7,\ U_{min}=\frac{3(4-\sqrt2)}{7}$.\\
	Задача решена.
\end{example}
\subsection{Дифференциальное исчисление в линейных нормированных пространствах}
\begin{definition}
	Линейный оператор $A:E_1\ra E_2$, где $E_1,E_2$ - линейные нормированные пространства называется \textit{ограниченным}, если:
	\[(\ex M\in\R)(\fa x\in E_1)\ \b[Ax]_{E_2}\le M\cdot\b[x]_{E_1}\]
	Наименьшая из всех таких констант называется \textit{нормой оператора}.
\end{definition}
\begin{theorem}
	Линейный оператор $A:E_1\ra E_2$ ограничен т. и т. т. к. он непрерывен на $E_1$.
\end{theorem}
\begin{proof}
	$\La$: Пусть $A$ - непрерывен на $E_1$, но неограничен, то есть:
	\[(\fa M\in \R)(\ex x\in E_1)\ \b[Ax]_{E_2}>M\cdot\b[x]_{E_1}\]
	Возьмём $\ds M:=n\in\N,\ x_n\in E_1: \b[Ax_n]_{E_2}>n\cdot\b[x_n]_{E_1},\ y_n:=\frac{x_n}{n\cdot\b[x_n]_{E_1}},\ \b[y_n]_{E_1}-\frac1n\us{n\ra\i}{\ra}0$.\\
	То есть $y_n\ra0$ в $E_1$, следовательно $Ay_n\us{n\ra\i}\ra0$ в $E_2$, так как $A$ - непрерывно.\\
	\[Ay_n=\frac1{n\b[x_n]_{E_1}}\cdot Ax_n\Ra\b[Ay_n]_{E_2}=\frac1{n\cdot\b[x_n]_{E_1}}\cdot\b[Ax_n]_{E_2}>1\]
	$Ay_n\ra0$ при $n\ra\i$, но в то же время $\b[Ay_n]>1$. Противоречие, в одну сторону доказали.\\
	$\Ra$: Пусть A - ограничен. Докажем, что A непрерывен в точке $x_0\in E_1$.\\
	Пусть $x_n\ra x_0$ при $n\ra\i$, что равносильно $\b[x_n-x_0]_{E_1}\ra0$ при $n\ra\i$. Тогда:
	\[\b[Ax_n-Ax_0]_{E_2}=\b[A(x_n-x_0)]_{E_2}\le M\cdot\b[x_n-x_0]_{E_1}\ra0\text{ при }n\ra\i\]
	Что и означает непрерывность.\\
	Получили требуемое.
\end{proof}
\begin{definition}
	Норма оператора $\b[A]$ обладает, по определению, свойствами:
	\begin{enumerate}
		\item $(\fa x\in E_1)\ \b[Ax]_{E_2}\le\b[A]\cdot\b[x]_{E_1}$.
		\item $(\fa\eps>0)(\ex x_\eps\in E_1)\ \b[Ax_\eps]_{E_2}>(\b[A]-\eps)\cdot\b[x_\eps]_{E_1}$.
	\end{enumerate}
\end{definition}
\begin{proposition}
	Если $A$ - линейный ограниченный оператор из линейного нормированного пространства $E_1$ в линейное нормированное пространство $E_2$, то:
	\[\b[A]=\sup_{\b[x]_{E_1}\le1}\b[Ax]_{E_2}=\sup_{x\ne0}\frac{\b[Ax]_{E_2}}{\b[x]_{E_1}}=\sup_{\b[x]_{E_1}=1}\b[Ax]_{E_2}\]
\end{proposition}
\begin{proof}
	\[(\fa x\in E_1,\ \b[x]_{E_1}\le1)\ \b[Ax]_{E_2}\le\b[A]\Ra\b[A]\text{ - верхняя грань }\{\b[Ax]_{E_2}\}\]
	По свойству нормы:
	\[(\fa\eps>0)(\ex x_\eps\in E_1)\ \b[Ax_\eps]_{E_2}>(\b[A]-\eps)\cdot\b[x_\eps]_{E_1}\]
	Пусть $\ds\xi_\eps:=\frac{x_\eps}{\b[x_\eps]_{E_1}}\Ra\b[\xi_\eps]_{E_1}=1$, при этом:
	\[\b[A\xi_\eps]_{E_2}>\b[A]-\eps\]
	Получили требуемое.
\end{proof}
\begin{note}
	Множество всех линейных ограниченных оператором из $E_1$ в $E_2$ является линейным нормированным пространством с нормой как раз той, которую только что определили, обозначается как $\mc L(E_1\ra E_2)$.
\end{note}
\begin{definition}
	Пусть $F$ - отображение открытого множества $G\subset E_1$, $E_1$ - л.н.п., в л.н.п. $E_2$. $F$ называется \textit{(сильно) дифференцируемым} в $x_0\in G$, если $\ex$ линейный ограниченный оператор $L_{x_0}:E_1\ra E_2$ такой, что:
	\[\b[F(x_0+h)-F(x_0)-L_{x_0}h]_{E_2}=\sigma(\b[h]_{E_1}),\ h\ra0\text{ в }E_1\] 
	$L_{x_0}$ называется \textit{(сильной) производной (производной Фреше)} отображения $F$ в точке $x_0$, $L_{x_0}=f'(x_0)$.
\end{definition}